\documentclass[a4paper,12pt]{article}

\usepackage[utf8]{inputenc}
\usepackage[russian]{babel}
\usepackage{geometry}
\geometry{margin=2.5cm}

\title{Мини-движок квеста на C\# \\ Проектирование архитектуры \\ Урок 1}
\author{}
\date{}

\begin{document}

\maketitle

\section{Список классов и интерфейсов}

\subsection{Основные классы}
\begin{itemize}
\item Game
\item GameState
\item Player
\item Location
\end{itemize}

\subsection{Интерфейсы}
\begin{itemize}
\item ICommand
\item IInteractable
\item ICondition
\item IEffect
\end{itemize}

\subsection{Абстрактные классы}
\begin{itemize}
\item CommandBase
\item ConditionBase
\item EffectBase
\item GameEventBase
\end{itemize}

\subsection{Команды}
\begin{itemize}
\item LookCommand
\item GoCommand
\item InteractCommand
\item InventoryCommand
\item HelpCommand
\end{itemize}

\subsection{Условия}
\begin{itemize}
\item HasItemCondition
\item FlagCondition
\item HealthCondition
\item AndCondition
\item OrCondition
\item NotCondition
\end{itemize}

\subsection{Эффекты}
\begin{itemize}
\item AddItemEffect
\item RemoveItemEffect
\item SetFlagEffect
\item DamageEffect
\item HealEffect
\item LogEffect
\item AddExitEffect
\item ChangeLocationEffect
\end{itemize}

\subsection{Объекты взаимодействия}
\begin{itemize}
\item Chest
\item Door
\item NPC
\item Trap
\end{itemize}

\subsection{События}
\begin{itemize}
\item OnEnterLocationEvent
\item OnTurnEvent
\item OneTimeEvent
\end{itemize}

\section{Роли основных классов}

\textbf{Game} — основной класс игры, управляет игровым циклом.

\textbf{GameState} — хранит состояние игры: здоровье игрока, инвентарь, флаги мира и журнал событий.

\textbf{Player} — представляет игрока и его характеристики.

\textbf{Location} — описывает локацию игры: название, описание, объекты, события и выходы.

\textbf{ICommand} — интерфейс всех игровых команд.

\textbf{IInteractable} — интерфейс объектов, с которыми может взаимодействовать игрок.

\textbf{ICondition} — интерфейс условий, проверяющих состояние игры.

\textbf{IEffect} — интерфейс эффектов, которые изменяют состояние игры.

\section{Схема локаций}

Игровой мир состоит из следующих локаций:

\begin{itemize}
\item Hall — стартовая зона
\item Storage — склад
\item DarkCorridor — тёмный коридор
\item GeneratorRoom — комната генератора
\item Exit — выход
\end{itemize}

Связи между локациями:

\begin{itemize}
\item Hall → Storage
\item Hall → DarkCorridor
\item Storage → GeneratorRoom
\item DarkCorridor → Exit
\end{itemize}

Некоторые переходы могут открываться только после выполнения условий.

\section{Игровые механики}

В игре реализованы следующие механики:

\begin{itemize}

\item \textbf{Тёмный коридор} — если игрок входит в DarkCorridor без предмета Torch, он получает урон.

\item \textbf{Дверь с ключом} — дверь в Storage открывается только если у игрока есть предмет Key.

\item \textbf{Генератор} — генератор можно включить только при наличии Fuse и Wrench.

\item \textbf{Ловушка} — в одной из локаций находится одноразовая ловушка, которая наносит урон игроку.

\end{itemize}

\section{Квесты}

\subsection{Квест 1: Включить генератор}

\textbf{Описание:}  
Игрок должен найти необходимые предметы и включить генератор.

\textbf{Условия выполнения:}

\begin{itemize}
\item Найти Fuse
\item Найти Wrench
\item Активировать генератор в GeneratorRoom
\end{itemize}

После выполнения устанавливается флаг \texttt{GeneratorOn}.

\subsection{Квест 2: Добраться до выхода}

\textbf{Описание:}  
Игрок должен найти путь через локации и добраться до Exit.

\textbf{Условие завершения:}

Игрок находится в локации Exit.

\end{document}